\begin{notechapter}
\chapter{The Perceptron}\label{ch:perceptron}
\chaptersubtitle{An update rule, its geometry, and why it stops}

\begin{chapterintro}
The perceptron (Rosenblatt, 1958) is the simplest learning algorithm for a linear
classifier: look at the points one at a time and, whenever one is misclassified,
add it to the weight vector (or subtract it). This chapter studies that rule from
three sides. Geometrically, each update rotates the decision boundary around a
point that can be computed exactly. Algorithmically, the rule is stochastic
(sub)gradient descent on a particular loss. Theoretically, it stops after a number
of mistakes bounded by the geometry of the data, when the data can be separated at
all. You will learn
\begin{itemize}[nosep]
  \item the update rule in augmented coordinates, and what it does to the score of
        the misclassified point;
  \item the pivot of the rotation, obtained as a cross product;
  \item the perceptron as gradient descent on $\max(0, -yf(x))$;
  \item Novikoff's theorem: at most $(R/\gamma)^2$ mistakes on separable data;
  \item the dual (kernel) perceptron, and what happens on non-separable data.
\end{itemize}
\prerequisites{linear algebra in $\R^2$ and $\R^3$; \cref{ch:surrogate-losses} for the margin.}
\end{chapterintro}

\section{The rule}\label{perceptron:sec:rule}

Labels are $y\in\{-1, +1\}$ (replace $y\in\{0,1\}$ by $2y - 1$). A linear classifier predicts
$\hat y = \sign(w^\top x + b)$, with $\sign(0) = +1$ by convention. It is convenient to absorb the
bias into the weights with \emph{augmented} vectors
\[
  \tilde x = \begin{pmatrix}x\\1\end{pmatrix}\in\R^{d+1},\qquad
  \theta = \begin{pmatrix}w\\b\end{pmatrix}\in\R^{d+1},\qquad
  w^\top x + b = \theta^\top\tilde x .
\]
A point is correctly classified when its \emph{functional margin} $y\,\theta^\top\tilde x$ is positive.

\begin{definition}[Perceptron]\label{perceptron:def:rule}
Start from $\theta_0$ (usually $0$). Visit the training points in some order, repeatedly.
When the visited point $(x, y)$ satisfies $y\,\theta^\top\tilde x\leq0$ (a \emph{mistake}), update
\[
  \theta\leftarrow\theta + \alpha\,y\,\tilde x,\qquad\text{i.e.}\qquad
  w\leftarrow w + \alpha\,y\,x,\quad b\leftarrow b + \alpha\,y,
\]
with a learning rate $\alpha>0$. Otherwise leave $\theta$ unchanged. Stop after a full pass without
mistakes.
\end{definition}

In two dimensions, with $x = (p, q)$, a misclassified positive point is added,
$(w_1, w_2, b)\leftarrow(w_1 + \alpha p, w_2 + \alpha q, b + \alpha)$, and a misclassified negative point is
subtracted. The effect on the point itself is immediate:
\[
  y\,(\theta + \alpha y\tilde x)^\top\tilde x = y\,\theta^\top\tilde x + \alpha\,y^2\norm{\tilde x}^2 = y\,\theta^\top\tilde x + \alpha\bigl(\norm x^2 + 1\bigr):
\]
its functional margin increases by $\alpha(\norm x^2 + 1)>0$. One update may not suffice to
classify it correctly, and it may spoil other points; the theorem of
\cref{perceptron:sec:convergence} shows that this does not go on forever.

\paragraph{The learning rate does not matter.} If $\theta_0 = 0$, every $\theta$ produced by the
algorithm is $\alpha$ times the one produced with $\alpha = 1$, and $\sign(\alpha\theta^\top\tilde x) = \sign(\theta^\top\tilde x)$:
the sequence of mistakes and the final classifier are the same for every $\alpha$. We
take $\alpha = 1$ from now on.

\section{The geometry of one update}\label{perceptron:sec:geometry}

The decision boundary is the line (hyperplane, for $d>2$) $\{x : w^\top x + b = 0\}$. Its
normal vector is $w$: for two points $x, x'$ on it, $w^\top(x - x') = 0$. The \emph{signed distance}
of a point to the boundary is
\[
  d(x) = \frac{w^\top x + b}{\norm w},
\]
positive on the side $w$ points to. The functional margin $y(w^\top x + b)$ is $\norm w$ times the
signed distance of $x$ to the boundary, counted positively on the correct side.

\paragraph{Rotation around a pivot.} After an update on $x = (p, q)$, the new boundary is
\[
  (w_1 + yp)x_1 + (w_2 + yq)x_2 + (b + y) = 0
  \iff \underbrace{w_1x_1 + w_2x_2 + b}_{\text{old boundary}} + y\,\underbrace{(px_1 + qx_2 + 1)}_{\text{line }L_x} = 0 .
\]
Every point on both the old boundary and the line $L_x : px_1 + qx_2 + 1 = 0$ satisfies the new
equation. If the two lines meet, the new boundary passes through their intersection
$x_0$: the update \emph{rotates the boundary around $x_0$}. The line $L_x$ depends only on the
misclassified point: its normal vector is $(p, q) = x$, so it is perpendicular to the
direction of $x$, and its signed distance to the origin is $-1/\norm x$. It lies on the other side
of the origin from $x$, at distance $1/\norm x$.

\paragraph{The pivot as a cross product.} Augmented coordinates make the
computation transparent. A point $x$ of the plane corresponds to the vector $\tilde x = (x_1, x_2, 1)$,
and the boundary $\{x : \theta^\top\tilde x = 0\}$ is the trace, on the plane $\{u_3 = 1\}$ of $\R^3$, of the plane
through the origin with normal $\theta$. Likewise, $L_x$ is the trace of the plane with normal $\tilde x$.
Two planes through the origin meet along the line spanned by the cross product of
their normals, so the pivot is that line's point with third coordinate $1$:
\[
  \theta\times\tilde x =
  \begin{pmatrix}w_1\\w_2\\b\end{pmatrix}\times\begin{pmatrix}p\\q\\1\end{pmatrix}
  = \begin{pmatrix}w_2 - bq\\bp - w_1\\w_1q - w_2p\end{pmatrix},
  \qquad
  x_0 = \frac1{w_1q - w_2p}\begin{pmatrix}w_2 - bq\\bp - w_1\end{pmatrix}.
\]
The new normal $\theta + y\tilde x$ is also orthogonal to $\theta\times\tilde x$, so the new boundary contains
$x_0$, as found above. When $w_1q - w_2p = 0$, that is, when $w$ is parallel to $x$, the two
lines are parallel and there is no pivot: the update translates the boundary
without rotating it (the pivot ``is at infinity'').

\begin{keyidea}[A pencil of lines]
All lines through $x_0$ have equations $\theta^\top\tilde x + t\,\tilde x_{\mathrm{mis}}^\top\tilde x = 0$, $t\in\R$, where $\tilde x_{\mathrm{mis}}$ is the
misclassified point: this family is the \emph{pencil} of lines through $x_0$. An update moves the
boundary from $t = 0$ to $t = \alpha y$ along this pencil. As $\alpha\to\infty$ the boundary tends to $L_x$; a
single update can therefore never carry the boundary past $L_x$.
\end{keyidea}

\begin{example}[Three updates]\label{perceptron:ex:three}
Take the six points of \cref{perceptron:fig:example},
\begin{center}
\begin{tabular}{@{}lcccccc@{}}
\toprule
$i$ & 1 & 2 & 3 & 4 & 5 & 6 \\
\midrule
$x_i$ & $(2, 1)$ & $(0, -1)$ & $(1, 2.5)$ & $(-1, 1)$ & $(2.5, 2.5)$ & $(-2, -1)$ \\
$y_i$ & $+1$ & $-1$ & $+1$ & $-1$ & $+1$ & $-1$ \\
\bottomrule
\end{tabular}
\end{center}
visited in this order from $\theta_0 = 0$, with $\alpha = 1$.
\begin{itemize}[nosep]
  \item Point 1: $\theta_0^\top\tilde x_1 = 0$, a mistake. $\theta_1 = (2, 1, 1)$: boundary $2x_1 + x_2 + 1 = 0$.
  \item Point 2: $y_2\theta_1^\top\tilde x_2 = -(0 - 1 + 1) = 0$, a mistake. $\theta_2 = \theta_1 - (0, -1, 1) = (2, 2, 0)$: boundary
        $x_1 + x_2 = 0$. The pivot is $\theta_1\times\tilde x_2 = (2, -2, -2)$, i.e.\ $x_0 = (-1, 1)$, which is indeed on both
        boundaries.
  \item Point 3: $\theta_2^\top\tilde x_3 = 7>0$, correct.
  \item Point 4: $y_4\theta_2^\top\tilde x_4 = -(-2 + 2) = 0$, a mistake. $\theta_3 = (2, 2, 0) - (-1, 1, 1) = (3, 1, -1)$: boundary
        $3x_1 + x_2 - 1 = 0$, pivot $\theta_2\times\tilde x_4 = (2, -2, 4)$, i.e.\ $x_0 = (0.5, -0.5)$.
  \item Points 5 and 6 are correct; a second pass finds no mistake.
\end{itemize}
The algorithm stops after three mistakes with $w = (3, 1)$, $b = -1$.
\end{example}

\begin{figure}[ht]
\centering
\begin{tikzpicture}
\begin{axis}[width=9.5cm, axis equal image, xmin=-3, xmax=3.5, ymin=-2.2, ymax=3.3,
  xlabel={$x_1$}, ylabel={$x_2$}, axis lines=middle, xtick={-2,...,3}, ytick={-2,...,3},
  tick label style={font=\scriptsize}, clip=true]
  \addplot[thick, gray, densely dotted, domain=-3:3.5, samples=2] {-2*x-1};
  \addplot[thick, bookexample, dashed, domain=-3:3.5, samples=2] {-x};
  \addplot[very thick, bookaccent, domain=-3:3.5, samples=2] {1-3*x};
  \addplot[only marks, mark=*, mark size=2.2pt, bookdef] coordinates {(2,1) (1,2.5) (2.5,2.5)};
  \addplot[only marks, mark=square*, mark size=2.2pt, bookpitfall] coordinates {(0,-1) (-1,1) (-2,-1)};
  \addplot[only marks, mark=o, mark size=3.5pt, black, thick] coordinates {(-1,1) (0.5,-0.5)};
  \node[font=\scriptsize, anchor=south east] at (axis cs:-1.05,1.05) {pivot 1};
  \node[font=\scriptsize, anchor=north west] at (axis cs:0.55,-0.55) {pivot 2};
  \node[font=\scriptsize, gray, anchor=west] at (axis cs:-2.9,2.3) {$\theta_1$};
  \node[font=\scriptsize, bookexample, anchor=west] at (axis cs:-2.9,3.0) {$\theta_2$};
  \node[font=\scriptsize, bookaccent, anchor=west] at (axis cs:-0.45,2.9) {$\theta_3$};
\end{axis}
\end{tikzpicture}
\caption{\Cref{perceptron:ex:three}: positive points (green discs), negative points (red
squares), and the boundaries after each update: $\theta_1$ (dotted), $\theta_2$ (dashed), $\theta_3$ (solid,
final). The second update rotates the boundary around $(-1, 1)$, the third around $(0.5, -0.5)$;
each pivot is the intersection of the old boundary with the line $L_x$ of the misclassified
point.}
\label{perceptron:fig:example}
\end{figure}

\section{The perceptron as gradient descent}\label{perceptron:sec:sgd}

With the margin $m = y\,\theta^\top\tilde x$ of \cref{ch:surrogate-losses}, consider the \emph{perceptron loss}
\[
  \ell(\theta; x, y) = \max\bigl(0, -y\,\theta^\top\tilde x\bigr) = \max(0, -m).
\]
It is zero for correctly classified points and grows linearly with the margin of a
misclassified one. Its gradient is $0$ when $m>0$ and $-y\tilde x$ when $m<0$; at $m = 0$, $-y\tilde x$ is a
subgradient. A stochastic subgradient step on one point is therefore
\[
  \theta\leftarrow\theta - \alpha\,\partial\ell = \theta + \alpha\,y\tilde x\quad\text{if }m\leq0,
\]
exactly the perceptron. Compared with the hinge loss $\max(0, 1 - m)$ of the support vector
machine (\cref{ch:svm}), the perceptron loss has its corner at $m = 0$ instead of $m = 1$: it is
satisfied by any separating hyperplane, however close to the points. This is why its
solution depends on the order of the points and the initialization, and why it has
no margin: the final boundary of \cref{perceptron:ex:three} passes at distance
$2/\sqrt{10}\approx0.63$ from $(0, -1)$, whereas the best separating line keeps every point at distance at
least $1.25$ (\cref{ch:svm}). The
perceptron loss is not classification-calibrated in the sense of
\cref{surrogate-losses:thm:bjm} ($\phi(m) = \max(0, -m)$ has $\phi'(0) = 0$ from the right): the zero
function minimizes it.

\section{Convergence on separable data}\label{perceptron:sec:convergence}

\begin{theorem}[Novikoff, 1962]\label{perceptron:thm:novikoff}
Suppose $\norm{\tilde x_i}\leq R$ for all $i$, and that some unit vector $u\in\R^{d+1}$ separates the data with
margin $\gamma>0$: $y_i\,u^\top\tilde x_i\geq\gamma$ for all $i$. Then the perceptron started at $\theta_0 = 0$ makes at most
$(R/\gamma)^2$ mistakes, whatever the order of the points.
\end{theorem}

\begin{proof}
Let $\theta_k$ be the weight vector after $k$ mistakes, the $k$-th on $(\tilde x, y)$.
\begin{itemize}[nosep]
  \item \emph{It moves towards $u$}: $u^\top\theta_k = u^\top\theta_{k-1} + y\,u^\top\tilde x\geq u^\top\theta_{k-1} + \gamma$, so $u^\top\theta_k\geq k\gamma$.
  \item \emph{It grows slowly}: $\norm{\theta_k}^2 = \norm{\theta_{k-1}}^2 + 2y\,\theta_{k-1}^\top\tilde x + \norm{\tilde x}^2\leq\norm{\theta_{k-1}}^2 + R^2$, because the
        middle term is $\leq0$ on a mistake. So $\norm{\theta_k}^2\leq kR^2$.
\end{itemize}
By the Cauchy--Schwarz inequality, $k\gamma\leq u^\top\theta_k\leq\norm{\theta_k}\leq\sqrt k\,R$, hence $k\leq R^2/\gamma^2$.
\end{proof}

The bound does not depend on the number of points or on the dimension, only on the
ratio of the size of the data to the margin. It is a worst case: in
\cref{perceptron:ex:three}, $R = \norm{(2.5, 2.5, 1)}\approx3.67$, and the best margin in augmented
coordinates is $\gamma\approx0.97$ (\cref{perceptron:exr:2}), so the bound is $14.4$, against three
actual mistakes. Note that $\gamma$ is measured in the augmented space $\R^{d+1}$, with the bias
treated as a weight; it is smaller than the geometric margin in $\R^d$ when the
separating line is far from the origin. Centering the data helps.

\paragraph{Non-separable data.} If no hyperplane separates the data, mistakes
continue forever and $\theta$ does not converge, although it stays bounded (the
\emph{perceptron cycling theorem}). Practical variants keep the best weights seen
so far (the \emph{pocket} algorithm), or average all the iterates, which also
improves generalization on separable data (the \emph{averaged} or \emph{voted}
perceptron of Freund and Schapire, 1999).

\section{The dual form and kernels}\label{perceptron:sec:dual}

Starting from $0$, $\theta$ is always a combination of the training points:
\[
  \theta = \sum_{i=1}^nc_i\,y_i\,\tilde x_i,\qquad c_i = \text{number of mistakes made on point }i .
\]
The score of a new point is then $\theta^\top\tilde x = \sum_ic_iy_i\,\tilde x_i^\top\tilde x$, and the update on point $j$ is
simply $c_j\leftarrow c_j + 1$. Only inner products between points are used. Replacing
$\tilde x_i^\top\tilde x$ by a kernel $k(x_i, x) + 1$ (\cref{svm:sec:kernels}) gives the \emph{kernel perceptron},
whose boundary $\sum_ic_iy_i(k(x_i, x) + 1) = 0$ can be nonlinear in $x$; Novikoff's theorem holds in the
feature space of the kernel. In \cref{perceptron:ex:three}, $c = (1, 1, 0, 1, 0, 0)$:
$\theta = \tilde x_1 - \tilde x_2 - \tilde x_4 = (2, 1, 1) - (0, -1, 1) - (-1, 1, 1) = (3, 1, -1)$.

\begin{chaptersummary}
\begin{itemize}[nosep]
  \item The perceptron adds $y\tilde x$ to $\theta$ on each mistake; this raises the functional margin
        of the point by $\norm x^2 + 1$. With $\theta_0 = 0$ the learning rate is irrelevant.
  \item An update rotates the boundary around the intersection $x_0$ of the old boundary
        with the line $L_x : x^\top z + 1 = 0$; in augmented coordinates $x_0$ is given by $\theta\times\tilde x$. If
        $w\parallel x$, the boundary is translated instead.
  \item The rule is stochastic subgradient descent on $\max(0, -y\theta^\top\tilde x)$, a loss without a margin.
  \item On data separable with margin $\gamma$ in a ball of radius $R$, it makes at most $(R/\gamma)^2$
        mistakes; on non-separable data it does not converge.
  \item $\theta$ is a combination of training points with mistake counts as coefficients, which
        gives the kernel perceptron.
\end{itemize}
\end{chaptersummary}

\section*{Exercises}
\addcontentsline{toc}{section}{Exercises}

\begin{exercise}\label{perceptron:exr:1}
Show that the line $L_x : px_1 + qx_2 + 1 = 0$ is perpendicular to $x = (p, q)$ and lies at distance
$1/\norm x$ from the origin, on the opposite side from $x$. Check the formula for the pivot by
verifying directly that $x_0$ satisfies both $w^\top x_0 + b = 0$ and $x^\top x_0 + 1 = 0$.
\solutionlink{perceptron:exr:1}
\end{exercise}

\begin{exercise}\label{perceptron:exr:2}
For the data of \cref{perceptron:ex:three}, show that $v = (8, 4, -7)/11$ satisfies $y_iv^\top\tilde x_i\geq1$ for all $i$,
with equality for three points. Deduce that $u = v/\norm v$ separates the augmented points with
margin at least $1/\norm v\approx0.97$. (One can show that no unit vector does better, so this is the $\gamma$ of
Novikoff's theorem.)
\solutionlink{perceptron:exr:2}
\end{exercise}

\begin{exercise}\label{perceptron:exr:3}
Run the perceptron on the data of \cref{perceptron:ex:three} visited in reverse order
(point 6 first). How many mistakes does it make, and which boundary does it return?
\solutionlink{perceptron:exr:3}
\end{exercise}

\begin{exercise}\label{perceptron:exr:4}
Show that the four points of the XOR problem, $(0,0)$ and $(1,1)$ with label $-1$, $(0,1)$ and $(1,0)$
with label $+1$, are not linearly separable. Show that they become separable with the
kernel $k(x, x') = (x^\top x')^2$, and give the corresponding feature map.
\solutionlink{perceptron:exr:4}
\end{exercise}
\end{notechapter}
